\chapter*{Abstract}
\label{chap:Abstract}

La tesi analizza e confronta due tecnologie per il rilevamento della presenza umana in ambienti indoor, il sensore PIR (Passive InfraRed) e il radar mmWave a 24 GHz, per stabilire quale delle due rilevi una persona con maggiore affidabilità in uno scenario post-sisma. Il contesto applicativo in cui il confronto trova riscontro è il progetto DIPME, che prevede l’integrazione di sensori IoT (Internet of Things) negli arredi scolastici per segnalare, in seguito a un evento sismico, l’eventuale presenza di persone rifugiate sotto di essi.

Il nodo di quel progetto affida oggi il rilevamento della presenza a un sensore PIR. Tuttavia, a causa del suo principio di funzionamento, il PIR rileva il movimento ma non una persona immobile, cioè proprio il caso che il progetto deve coprire. Il lavoro verifica quindi se il radar, di tipo FMCW (Frequency Modulated Continuous Wave), possa superare questo limite e con quali compromessi.

Nella prima parte vengono studiati il sensore PIR \pir{} e i due moduli radar \ldb{} e \ldc{}, con particolare attenzione ai loro principi di funzionamento, ai protocolli di comunicazione, ai parametri disponibili e ai limiti riportati nella documentazione dei costruttori. Per il confronto è stata realizzata una piattaforma di acquisizione basata su ESP32, in grado di campionare i sensori in modo sincrono e di accedere, nel caso dei radar, all’energia riflessa per ciascuna cella di distanza. A supporto della campagna è stata progettata e implementata un’interfaccia web, servita direttamente dall’ESP32, per la visualizzazione dei dati in tempo reale e la loro esportazione nello stesso formato usato dalla piattaforma di acquisizione.

La piattaforma è stata utilizzata per una campagna sperimentale di 444 trial che ha preso in esame la persona ferma e in movimento, la distanza di rilevamento, le latenze, i falsi positivi, l’attenuazione del segnale in presenza di ostacoli, la copertura angolare e il rilevamento del respiro. I risultati confermano il limite del PIR. Con la persona immobile, anche sotto il banco, il PIR l’ha rilevata in meno del 2\,\% del tempo, mentre il radar in tutti gli scenari di confronto. Anche il secondo radar ha confermato questo comportamento, con limiti che la tesi attribuisce all’esemplare in prova.

Sulle grandezze graduate fornite dal radar è stato inoltre sviluppato un indice di vitalità calcolato direttamente a bordo, che, oltre alla presenza, stima quanto la persona si muove. Il consumo energetico dei sensori è stato infine valutato sulla base dei datasheet. Poiché il radar consuma circa mille volte più del PIR, la tesi conclude che il radar deve affiancare il PIR e non sostituirlo, restando spento nel funzionamento ordinario e acceso quando il nodo entra nel regime di emergenza.
